\documentclass[11pt,a4paper]{article}

\usepackage[a4paper,margin=1in]{geometry}
\usepackage{fontspec}
\IfFontExistsTF{Times New Roman}
  {\setmainfont{Times New Roman}}
  {\IfFontExistsTF{TeX Gyre Termes}{\setmainfont{TeX Gyre Termes}}{\setmainfont{Times}}}
\usepackage{newtxmath}
\usepackage{amsmath}
\usepackage{booktabs}
\usepackage{longtable}
\usepackage{array}
\usepackage{enumitem}
\usepackage{xcolor}
\usepackage{listings}
\usepackage[hidelinks]{hyperref}
\usepackage{xurl}
\usepackage{fancyhdr}
\usepackage{microtype}

\setlength{\parindent}{0pt}
\setlength{\parskip}{0.65em}
\setlist[itemize]{leftmargin=1.4em,itemsep=0.2em,topsep=0.2em}
\setlist[enumerate]{leftmargin=1.6em,itemsep=0.2em,topsep=0.2em}

\pagestyle{fancy}
\fancyhf{}
\lhead{\small Ticket 02 Study Report}
\rhead{\small Grid and Operator Setup}
\cfoot{\thepage}
\renewcommand{\headrulewidth}{0.3pt}

\lstset{
  basicstyle=\ttfamily\small,
  breaklines=true,
  frame=single,
  columns=fullflexible,
  keepspaces=true,
  rulecolor=\color{black!30},
  xleftmargin=0.5em,
  xrightmargin=0.5em
}

\newcommand{\code}[1]{\texttt{#1}}
\newcommand{\file}[1]{\path{#1}}
\newcolumntype{L}[1]{>{\raggedright\arraybackslash}p{#1}}
\newcolumntype{P}[1]{>{\raggedright\arraybackslash\ttfamily\footnotesize}p{#1}}

\title{\textbf{Ticket 02 Study Report: Finite-Difference Grid and Black-Scholes Operator Setup}\\
\large SURF2026 American Option Risk Surfaces}
\author{Codex-assisted implementation report}
\date{June 15, 2026}

\begin{document}
\maketitle

\begin{abstract}
Ticket 2 builds the first numerical layer beneath the future option solver: uniform spot and
time-to-maturity grids, an explicit interior-node convention, central finite-difference
coefficients for the Black-Scholes spatial operator, a small operator-application helper, and
simple boundary-value helpers. This report is written for beginner researchers with accounting,
data, machine-learning, and Python experience who are still learning finite differences,
PDE-based pricing, CN/PSOR, and American option free-boundary logic. Ticket 2 prepares later
European Crank-Nicolson validation and American CN/PSOR work, but it does not implement or
validate those solvers.
\end{abstract}

\tableofcontents
\newpage

\section{Role of Ticket 2 in the Whole Solver Project}

Ticket 1 gave the project analytic option-pricing utilities: payoff functions and European
Black-Scholes closed-form prices. Ticket 2 begins the numerical side. Its purpose is to define
the finite grid on which a future solver will approximate option values, and to define the
discrete spatial operator that will later appear inside the Crank-Nicolson matrices.

This order matters because CN/PSOR combines several ideas at once: a grid, finite differences,
time stepping, linear algebra, an American payoff obstacle, and iterative projection. If all of
those pieces are written together, debugging becomes confusing. Ticket 2 isolates only the grid
and the spatial operator. Later tickets can then test time stepping and projection against a
known, reviewed numerical foundation.

In plain language, Ticket 2 answers three basic questions:
\begin{itemize}
  \item Where are the spot and time points?
  \item Which spot nodes are solved by the PDE operator, and which nodes are imposed boundaries?
  \item What are the lower, diagonal, and upper finite-difference weights for the Black-Scholes
        spatial operator?
\end{itemize}

It is important to state the boundary of this ticket. Ticket 2 does not march through time, does
not solve a linear system, does not run PSOR, does not enforce the American obstacle, and does
not extract a free boundary. It creates reusable pieces that Ticket 3 can use for the European
Crank-Nicolson validation check.

\section{Theory and Methodology Connection}

The source basis for this ticket is internal and project-specific:
\begin{itemize}
  \item \file{reports/01_literature/literature_map.md};
  \item \file{reports/02_math/formulation_note.md};
  \item \file{reports/03_solver/solver_validation_plan.md};
  \item \file{reports/03_solver/solver_implementation_tickets.md};
  \item the Ticket 1 Black-Scholes utility report and code.
\end{itemize}

The literature map places this work in the classical finite-difference American option
literature. It mentions Brennan-Schwartz finite-difference valuation as background and emphasizes
that a trusted classical benchmark must come before neural surrogate modelling. This Ticket 2
report does not add a standalone bibliography because the available project notes do not provide
complete bibliographic details for every external source. Instead, it uses the internal literature
map and formulation note as the documented source basis.

The formulation note defines the dividend-adjusted Black-Scholes spatial operator:
\begin{equation}
  L U =
  \frac{1}{2}\sigma^2 S^2 U_{SS}
  + (r-q)S U_S
  - r U.
  \label{eq:continuous-operator}
\end{equation}

The solver validation plan states that the finite-difference version \(L_h\) should approximate
\(U_S\) and \(U_{SS}\) on interior nodes, while boundary conditions must be handled consistently
later. Ticket 2 implements this \(L_h\) coefficient setup only. The future CN step will use
\begin{equation}
  A = I - \frac{\Delta \tau}{2}L_h,
  \qquad
  b^n = \left(I + \frac{\Delta \tau}{2}L_h\right)U^n,
\end{equation}
but Ticket 2 deliberately does not build \(A\), \(b^n\), or any time-marching solver.

\section{Finite-Difference Grids in Plain Language}

A continuous option-pricing problem has infinitely many possible spot prices and times. A
computer cannot solve infinitely many points directly. A finite-difference method replaces the
continuous domain with a finite list of grid points.

For spot, Ticket 2 uses a uniform grid on a truncated domain:
\begin{equation}
  S_i = i\Delta S,\qquad i = 0,1,\ldots,M,
  \qquad
  \Delta S = \frac{S_{\max}}{M}.
\end{equation}

The domain starts at \(S=0\) and ends at \(S=S_{\max}\). The true Black-Scholes model allows
unbounded spot prices, but a numerical solver must choose a finite upper boundary. Later
validation work will check whether \(S_{\max}\) is large enough.

For time-to-maturity, Ticket 2 uses
\begin{equation}
  \tau_n = n\Delta \tau,\qquad n = 0,1,\ldots,N,
  \qquad
  \Delta \tau = \frac{T}{N}.
\end{equation}

The project uses \(\tau\), time-to-maturity, rather than calendar time. At \(\tau=0\), the option
is at maturity and the value is the payoff. As \(\tau\) increases, the future solver moves away
from maturity toward earlier calendar times.

Uniform grids are not the most advanced choice, but they are the right first step here. They make
indexing readable, make central finite-difference formulas simple, and keep the beginner
implementation auditable. Nonuniform grids can be considered later only after the uniform-grid
solver is validated.

\section{Black-Scholes Spatial Operator}

The operator in Equation~\ref{eq:continuous-operator} has three parts:
\begin{longtable}{L{0.30\textwidth}L{0.58\textwidth}}
\toprule
Term & Plain-language meaning \\
\midrule
\endhead
\(\frac{1}{2}\sigma^2S^2U_{SS}\) &
Diffusion or curvature term. It captures how volatility and option curvature affect value. \\
\((r-q)SU_S\) &
Drift term. It captures the risk-neutral movement of the underlying after continuous dividends. \\
\(-rU\) &
Discounting term. It subtracts the risk-free discounting effect from the value. \\
\bottomrule
\end{longtable}

Ticket 2 converts this continuous operator into a tridiagonal discrete operator. ``Tridiagonal''
means that at an interior node \(i\), the approximation uses only three neighboring values:
\(U_{i-1}\), \(U_i\), and \(U_{i+1}\).

\section{Central Finite Differences}

At an interior node \(S_i\), central differences approximate the first and second spot
derivatives:
\begin{align}
  U_S(S_i) &\approx \frac{U_{i+1} - U_{i-1}}{2\Delta S}, \\
  U_{SS}(S_i) &\approx \frac{U_{i+1} - 2U_i + U_{i-1}}{(\Delta S)^2}.
\end{align}

Substituting these formulas into Equation~\ref{eq:continuous-operator} gives
\begin{equation}
  (L_h U)_i = a_i U_{i-1} + b_i U_i + c_i U_{i+1},
\end{equation}
where
\begin{align}
  a_i &=
    \frac{1}{2}\sigma^2\frac{S_i^2}{(\Delta S)^2}
    - \frac{1}{2}(r-q)\frac{S_i}{\Delta S}, \\
  b_i &=
    -\sigma^2\frac{S_i^2}{(\Delta S)^2}
    - r, \\
  c_i &=
    \frac{1}{2}\sigma^2\frac{S_i^2}{(\Delta S)^2}
    + \frac{1}{2}(r-q)\frac{S_i}{\Delta S}.
\end{align}

The implementation calls these arrays \code{lower}, \code{diagonal}, and \code{upper}. The names
come from the matrix form: \(a_i\) multiplies the value to the left, \(b_i\) multiplies the
current node, and \(c_i\) multiplies the value to the right.

Two manufactured-function checks are useful:
\begin{align}
  L(1) &= -r, \\
  L(S) &= -qS.
\end{align}
These identities are simple but powerful. They check the signs of the drift and discount terms
without needing a full option solver.

\section{Interior Nodes and Boundary Nodes}

For a grid with nodes \(0,\ldots,M\), Ticket 2 defines:
\begin{itemize}
  \item node \(0\) as the lower boundary at \(S=0\);
  \item node \(M\) as the upper boundary at \(S=S_{\max}\);
  \item nodes \(1,\ldots,M-1\) as interior nodes.
\end{itemize}

The function \code{interior\_indices(M)} returns exactly those interior indices. This is a small
helper, but it prevents a common off-by-one error. The future PDE solve should operate on
interior nodes while boundary values are imposed separately.

Boundary values are included in the full value vector because the discrete operator at the first
and last interior nodes needs \(U_0\) and \(U_M\). However, the operator helper returns only
interior results. This keeps the distinction between imposed boundary values and solved interior
values explicit.

\section{Design Decisions}

\subsection{Uniform Grid First}

The implementation uses uniform spot and time grids because the formulation note and solver
validation plan introduce the method with \(\Delta S\) and \(\Delta \tau\). Uniform grids make
the coefficient formulas easy to inspect and test. This is especially important for greenhand
researchers who are still learning how the PDE becomes code.

\subsection{Explicit Interior Indexing}

Interior indexing is implemented as a named helper instead of being hidden inside slices
throughout the solver code. The goal is not abstraction for its own sake. The goal is to make it
obvious that nodes \(0\) and \(M\) are boundaries while nodes \(1,\ldots,M-1\) are the PDE
nodes.

\subsection{Coefficient Arrays Instead of Full Matrix}

Ticket 2 returns lower, diagonal, and upper coefficient arrays rather than a dense matrix. This is
enough for tridiagonal finite-difference work and avoids introducing matrix assembly details
before Ticket 3. The future CN solver can decide how to build or solve tridiagonal systems.

\subsection{Simple Boundary Helpers}

The boundary helpers return boundary values only. They do not adjust a Crank-Nicolson
right-hand side, solve a PDE, or extract a free boundary. The helpers are included because later
European validation and American cases need consistent boundary formulas.

\section{Implementation Record}

\subsection{Files Created}

\begin{longtable}{P{0.42\textwidth}L{0.50\textwidth}}
\toprule
File & Role in Ticket 2 \\
\midrule
\endhead
\file{src/american_risk_surfaces/solvers/grid.py} &
Ticket 02 module containing uniform spot/time grids plus explicit interior-node indices. \\
\file{src/american_risk_surfaces/solvers/operator.py} &
Ticket 02 module containing Black-Scholes operator coefficients, operator application, and boundary helpers. \\
\file{tests/test_grid_operator.py} &
Focused \code{unittest} coverage for grids, coefficients, manufactured-function checks, and boundaries. \\
\file{reports/03_solver/tickets/ticket_02_grid_operator.tex} &
Formal LaTeX study-report source. \\
\file{reports/03_solver/tickets/ticket_02_grid_operator.pdf} &
Formal PDF compiled from the LaTeX source. \\
\bottomrule
\end{longtable}

\section{Code-Level Function Explanations}

\begin{longtable}{P{0.39\textwidth}L{0.53\textwidth}}
\toprule
Function or class & Purpose and important behavior \\
\midrule
\endhead
\code{uniform\_spot\_grid} &
Builds \(S_i=i\Delta S\) for \(i=0,\ldots,M\). It validates \(S_{\max}>0\) and \(M\geq2\). \\
\code{uniform\_tau\_grid} &
Builds \(\tau_n=n\Delta\tau\) for \(n=0,\ldots,N\). It validates \(T\geq0\) and \(N\geq1\). \(T=0\) is allowed. \\
\code{interior\_indices} &
Returns \(1,\ldots,M-1\), making the boundary/interior convention explicit. \\
\code{BlackScholesOperatorCoefficients} &
Frozen dataclass holding lower, diagonal, upper, and interior spot arrays. \\
\code{black\_scholes\_operator\_coefficients} &
Computes \(a_i\), \(b_i\), and \(c_i\) on interior nodes using central differences. It rejects non-increasing, nonuniform, or \(dS\)-mismatched spot grids. \\
\code{apply\_black\_scholes\_operator} &
Applies the tridiagonal operator to a full value vector and returns interior values only. \\
\code{european\_put\_boundaries} &
Returns \((K e^{-r\tau},0)\), the simple European put boundary pair. \\
\code{american\_put\_boundaries} &
Returns \((K,0)\), the simple American put boundary pair. \\
\code{european\_call\_boundaries} &
Returns \((0,\max(S_{\max}e^{-q\tau}-Ke^{-r\tau},0))\). \\
\code{american\_call\_boundaries} &
Returns \((0,\max(0,S_{\max}-K,S_{\max}e^{-q\tau}-Ke^{-r\tau}))\). \\
\bottomrule
\end{longtable}

All helpers validate the domain inputs needed for this ticket: positive strike, nonnegative
\(\tau\), positive \(S_{\max}\), nonnegative volatility, positive \(\Delta S\), and one-dimensional
spot/value vectors where required. The operator helper additionally validates that the spot grid
is strictly increasing and uniformly spaced consistently with the supplied \(\Delta S\).

\section{Testing Strategy}

The test suite is designed to prove the local behavior of Ticket 2, not the future solver. The
main categories are:
\begin{longtable}{L{0.25\textwidth}L{0.31\textwidth}L{0.32\textwidth}}
\toprule
Test category & What it proves & What it does not prove \\
\midrule
\endhead
Spot grid tests &
Endpoint, spacing, and length are correct. &
Whether \(S_{\max}\) is large enough for all option cases. \\
Tau grid tests &
Endpoint, spacing, length, and \(T=0\) behavior are correct. &
Any time-stepping accuracy. \\
Invalid grid tests &
Bad grid sizes are rejected clearly. &
A complete production input-validation policy. \\
Operator grid validation tests &
Non-increasing grids, nonuniform grids, and \(dS\)-mismatched grids are rejected. &
Any future nonuniform-grid method. \\
Interior-index tests &
Boundary nodes are excluded from the PDE interior. &
Correct future matrix assembly by itself. \\
Coefficient shape tests &
Coefficient arrays align with interior nodes. &
Correct time-step matrices. \\
Manual coefficient test &
The lower, diagonal, and upper formulas match a hand calculation. &
Full PDE-solver convergence. \\
Manufactured-function tests &
\(L_h(1)=-r\) and \(L_h(S)=-qS\) hold on the grid. &
Option-pricing accuracy for kinked payoffs. \\
Boundary helper tests &
Simple put/call boundary formulas are implemented consistently. &
Boundary right-hand-side adjustment in CN. \\
\bottomrule
\end{longtable}

\section{Test Results and Interpretation}

The exact test command was:
\begin{lstlisting}[language=bash]
PYTHONPYCACHEPREFIX=/private/tmp/codex_pycache PYTHONPATH=src python3 -m unittest discover -s tests
\end{lstlisting}

Final observed result:
\begin{lstlisting}
Ran 25 tests in 0.002s

OK
\end{lstlisting}

This means the existing Ticket 1 tests and the new Ticket 2 grid/operator tests passed in the
current Python runtime. It gives evidence that the local grid helpers, coefficient formulas,
operator application, stronger operator-grid validation, and boundary-value helpers behave as
expected for the tested cases.

It does not mean the CN solver is correct, because there is no CN solver yet. It also does not
mean the American obstacle logic is correct, because there is no PSOR or American projection yet.
The correct interpretation is narrow: Ticket 2 has a tested numerical setup layer for later
solver validation.

\section{Limitations}

Ticket 2 has deliberate limits:
\begin{itemize}
  \item no Crank-Nicolson time stepping;
  \item no PSOR;
  \item no American option solve;
  \item no linear complementarity solve;
  \item no boundary extraction;
  \item no Greeks;
  \item no experiments or generated result tables;
  \item no datasets, stress maps, or neural-network files;
  \item no proof that a chosen \(S_{\max}\), \(M\), or \(N\) is sufficient for production analysis.
\end{itemize}

The boundary helpers are only formulas for boundary values. They are not the complete boundary
adjustments that a CN right-hand side will need. That distinction should remain visible in Ticket
3.

\section{Lessons Learned / Study Notes}

The first lesson is that a finite-difference method is mostly bookkeeping before it is linear
algebra. If the grid and interior indices are wrong, a later solver can look sophisticated while
still solving the wrong discrete problem.

The second lesson is that the Black-Scholes operator has financial meaning. Diffusion comes from
volatility, drift comes from \(r-q\), and discounting comes from \(-rU\). The signs matter.
A wrong sign in the drift term can reverse dividend behavior and damage later American call
validation.

The third lesson is that boundary nodes are not ordinary solved nodes. They are imposed values
that help the interior equations near the edge of the truncated domain. Later solver reports
should continue to distinguish boundary behavior from interior PDE behavior.

The fourth lesson is that simple manufactured functions are useful. Checking \(L_h(1)=-r\) and
\(L_h(S)=-qS\) is easier to understand than a full option price, but it still catches important
operator mistakes.

The fifth lesson is to keep validation layered. Ticket 2 passing is good news for the grid and
operator layer. It is not evidence that CN/PSOR, free-boundary extraction, or Greeks are correct.

\section{Link to Ticket 3}

Ticket 3 should build the European Crank-Nicolson validation mode. It can reuse:
\begin{itemize}
  \item \code{uniform\_spot\_grid} and \code{uniform\_tau\_grid} for the numerical domain;
  \item \code{interior\_indices} for the solved PDE nodes;
  \item \code{black\_scholes\_operator\_coefficients} for the spatial operator;
  \item boundary helpers for European call and put validation cases;
  \item Ticket 1 closed-form prices as benchmark values.
\end{itemize}

Ticket 3 should compare finite-difference European prices against closed-form Black-Scholes
prices before any American projection is added. It still should not implement PSOR, American
option solving, boundary extraction, Greeks, experiments beyond the reviewed validation scope,
datasets, stress maps, or neural networks unless a later approved ticket explicitly asks for them.

\section{Reproducibility Record}

\subsection{Files Changed}

\begin{longtable}{P{0.48\textwidth}L{0.44\textwidth}}
\toprule
Path & Status in this report task \\
\midrule
\endhead
\file{src/american_risk_surfaces/solvers/grid.py} &
Created for grid and interior-index helpers; module docstring identifies Ticket 02. \\
\file{src/american_risk_surfaces/solvers/operator.py} &
Created for operator coefficients and boundary helpers; module docstring identifies Ticket 02. \\
\file{tests/test_grid_operator.py} &
Created for Ticket 2 \code{unittest} coverage. \\
\file{reports/03_solver/tickets/ticket_02_grid_operator.tex} &
Created as the formal LaTeX report source. \\
\file{reports/03_solver/tickets/ticket_02_grid_operator.pdf} &
Compiled from the LaTeX report source. \\
\bottomrule
\end{longtable}

\subsection{Commands}

Test command:
\begin{lstlisting}[language=bash]
PYTHONPYCACHEPREFIX=/private/tmp/codex_pycache PYTHONPATH=src python3 -m unittest discover -s tests
\end{lstlisting}

Compile command for Python syntax:
\begin{lstlisting}[language=bash]
PYTHONPYCACHEPREFIX=/private/tmp/codex_pycache PYTHONPATH=src python3 -m compileall -q src tests
\end{lstlisting}

LaTeX compile command:
\begin{lstlisting}[language=bash]
HOME=/private/tmp FC_CACHEDIR=/private/tmp/fontconfig-cache latexmk -xelatex -interaction=nonstopmode -halt-on-error -outdir=/private/tmp/ticket02_latex reports/03_solver/tickets/ticket_02_grid_operator.tex
\end{lstlisting}

PDF verification used \code{pdfinfo} and rendered page images with \code{pdftoppm}. The PDF is a
LaTeX-generated formal report on A4 paper with one-inch margins, page numbers, tables, monospace
commands and paths, and LaTeX math.

\end{document}
